\subsection{Exp 4: timing, latency and matched transport controls}\label{sec:live-timing}\label{sec:r7-main}
All 690 saved action trajectories execute through a live local tool server under six conditions, producing 4,140 new TLS sessions. Each real response gates the next action; real and shadow calls have separate states. The original 414/276 customer-disjoint split is retained, with attacks fitted and frozen before test capture. Lengths, directions and arrival times measure the tool RPC chain.

A complete tunnel sends 4,096-byte cells in both directions at 2 or 5\,ms target intervals; dummy cells execute no tools. A training-sized 133-slot horizon continues after the last response. Serial and two-in-flight shadow schedules use the frozen prefix-only count policy with constant per-call framing. The full union covers size/order, timing, combined features and learned flush removal; the common seven, including a 1D-CNN and forest-leaf neighbors, permit a matched comparison with offline WF adaptations.
\begin{table*}[t]\centering\footnotesize
\caption{Timing-aware Exp 4 comparison: standard prompt, 138 held-out episodes, 24 customers (5 cancel-positive, 7 return-positive). Each BA pair is cancel/return (\%). The common seven models permit a matched-suite comparison; the full union also retains all earlier attacks.}\label{tab:r7-controls}
\begin{tabular}{@{}llrrrrrr@{}}\toprule
View & Condition & Common-seven BA & Full-union BA & Upper BA & Bytes$\times$ & Calls$\times$ & RPC p95 (ms)\\\midrule
M & Undefended & 84.3/76.9 & 91.7/86.4 & 100.0/97.2 & 1.00 & 1.00 & 15.4\\
M & Per-call constant & 53.2/59.1 & 68.6/62.7 & 92.8/79.5 & 3.92 & 1.00 & 16.6\\
M & Serial shadow & 50.0/52.9 & 61.4/53.4 & 71.4/65.8 & 15.11 & 3.83 & 26.2\\
M & Two-in-flight shadow & 50.0/57.3 & 51.7/57.7 & 69.1/69.7 & 15.11 & 3.83 & 20.8\\
M & Full tunnel, 2 ms & 59.9/55.2 & 62.5/55.2 & 70.2/62.0 & 125.71 & 1.00 & 56.4\\
M & Full tunnel, 5 ms & 54.6/53.7 & 54.6/57.3 & 72.9/64.2 & 125.71 & 1.00 & 112.3\\
\midrule
S & Record buckets & 53.1/69.8 & 53.1/69.8 & 78.5/86.9 & 2.25 & 1.00 & --\\
S & Tamaraw 40/12 ms & 55.3/55.3 & 55.3/55.3 & 68.7/64.4 & 12.15 & 1.00 & --\\
S & Tamaraw 4/1.2 ms & 57.8/51.8 & 57.8/51.8 & 71.8/61.8 & 11.99 & 1.00 & --\\
S & FRONT & 59.0/69.3 & 59.0/69.3 & 69.9/81.8 & 2.51 & 1.00 & --\\
S & WTF-PAD adaptation & 51.9/62.7 & 51.9/62.7 & 64.1/79.6 & 1.88 & 1.00 & --\\
S & WT random-donor & 50.0/56.7 & 50.0/56.7 & 67.0/67.8 & 2.04 & 1.00 & --\\
S & FRONT, short calibration & 52.4/56.0 & 52.4/56.0 & 61.3/71.1 & 2.33 & 1.00 & --\\
\bottomrule\end{tabular}
\par\smallskip\raggedright M: measured local tool RPC on saved actions; S: offline traffic transformation. RPC p95 ends at the last real response, excluding LLM generation. S-row simulated spans are not RPC latency. Upper BA is the full-union 95\% bound; every condition exceeds 60\%. M unions contain 17--22 attacks; S unions contain the common seven. Ratios use prompt-specific undefended captures.
\end{table*}


\paragraph{Timing reveals distinguishable cover}
For standard-prompt serial shadow, adding timing/flush attacks to the size/order attacks increases cancellation BA from 50.0\% to 61.4\% (paired increase 11.4 points, 95\% interval [3.7,19.6]); return changes from 50.7\% to 53.4\%. Flush-phase detection reaches 90.4/92.3\% for serial/parallel schedules over 15,278 records from 276 test episodes. A separate cover-versus-real request classifier reaches 82.5/96.1\% on 7,639 requests.

We then compose the detector with intent inference: a training-fitted detector strips the predicted suffix before classification. Table~\ref{tab:r8-flush} shows the recovered category signal, with a strongest return point of 60.7\%. The full attack union retains both stripped and unstripped traces. This separates the ability to recognize padding phases from the ability to recover a task's operation.
\begin{table}[t]\centering\footnotesize\setlength{\tabcolsep}{3pt}
\caption{Exp 4 intent BA after learned flush removal (95\% customer-cluster interval). Each row: 138 test episodes, 24 customers. Maximum of three classifiers plus majority; the full attack union additionally keeps unstripped observations.}\label{tab:r8-flush}
\begin{tabular}{@{}llrr@{}}\toprule
Prompt & Schedule & Cancel & Return\\\midrule
Standard & Serial & 57.8 [50.0, 70.7] & 53.4 [50.0, 64.6]\\
Standard & Parallel & 50.0 [50.0, 64.0] & 50.0 [50.0, 56.7]\\
Rules & Serial & 50.1 [50.0, 55.6] & 60.7 [50.8, 70.9]\\
Rules & Parallel & 50.8 [50.0, 61.7] & 54.0 [50.0, 67.0]\\
\bottomrule\end{tabular}\end{table}


\paragraph{Match cost to the protected boundary}
Within measured RPC conditions (M), per-call constant framing uses 3.92$\times$ bytes and 16.6\,ms p95, versus serial shadow's 15.11$\times$ and 26.2\,ms; its full-union BA is higher (68.6/62.7\% versus 61.4/53.4\%), a measured cost--inference trade-off. Two-in-flight shadow lowers p95 to 20.8\,ms, versus 56.4/112.3\,ms for 2/5\,ms tunnels. Sessions preserve real responses and final live state; fixed tunnel sizes/counts still expose host-scheduled direction/order. Table~\ref{tab:r7-controls} keeps measured RPC times separate from offline schedule transformations.

Under the common seven attacks, offline random-donor Walkie-Talkie (WT, S) reaches 50.0/56.7\% BA at 2.04$\times$ bytes; measured serial shadow (M) reaches 50.0/52.9\% at 15.11$\times$. This descriptive, shared-metadata comparison identifies a lower-byte shaping candidate. Establishing its online trade-off requires RPC measurements and the extended timing/flush suite; the offline point is neither a theoretical byte lower bound nor demonstrated deployment superiority.

The WF adapters use the same saved traffic and training customers. WT samples donor traces without symmetric site pairing; WTF-PAD samples training inter-arrival histograms; Tamaraw and FRONT supply fixed-rate and front-loaded controls. Input-derived fragments preserve sizes, directions and order, enabling a controlled metadata comparison. Their simulated spans and measured RPC latency remain separate cost objects.

\paragraph{Emulated WAN delay and jitter}
The relay re-runs all 276 test episodes with 20\,ms/5\,ms and 50\,ms/10\,ms one-way delay/jitter. The frozen loopback flush detector transfers with 89.4\% BA (versus 90.4\%), so the cover phase stays recognizable. Jitter does mask micro-timing: serial-shadow full-union cancellation BA falls from 61.4\% to 50.0\% and 54.0\%. This masking depends on observer position; an observer adjacent to the tool server sees loopback-like timing. Fixed horizons fail under delay: the 2\,ms tunnel's 133-slot (266\,ms) horizon completes only 85/276 and 5/276 sessions, whereas shadow cover and per-call framing complete all sessions at 855.4/1,999.5\,ms serial-shadow p95. A tunnel horizon must therefore be sized to path RTT, at a corresponding byte cost.

\textbf{Finding.} The same-capture timing contrast establishes schedule leakage missed by length-only evaluation. Exp~3's 50.0/50.0\% and Exp~4's 61.4/53.4\% also differ in attacker and collection scope. Within measured conditions, cheaper per-call framing leaves higher full-union inference; offline WF adapters identify lower-byte candidates under the common seven. Selecting between transport shaping and dispatch isolation therefore requires both the receiver boundary and the matched cost object.
